\section{Grids, preorder windows, and stable centers}\label{sec:windows}

We now prepare the index windows used to find a hit at a normalized scale
$t\in[1,2]$.  Each edge of a finite ordered tree will carry one window.
We will bound the set of centers at which translations by $tq^n$ from a
given window can change earlier grid keys, and show that the translated
points take distinct keys at the edge's own resolution.
Throughout this section, $q\in(0,1)$ and the integers
$M\ge2$, $d\ge1$, $g\ge1$, and $r_0\ge1$ are arbitrary and fixed.

For each integer $b\ge1$, put
\[
 c=\frac{4}{1-q},\qquad
 N_b=2^{\lceil\log_2(cq^{-b})\rceil},\qquad
 \Gamma_b=N_b^{-1}\Z.
\]
The periodic grid key of $z\in\R$ is
$J_b(z)=\lfloor N_b\{z\}\rfloor$, where
$\{z\}=z-\lfloor z\rfloor$.  Thus the cells are closed on the left and
open on the right, and a boundary point belongs to the cell on its right.
We have
\begin{equation}\label{eq:grid-size}
 cq^{-b}\le N_b<2cq^{-b}.
\end{equation}
All $N_b$ are powers of two and are nondecreasing in $b$.  Consequently,
if $b'\ge b$, then $N_{b'}/N_b$ is an integer and
\[
 J_b(z)=\left\lfloor\frac{J_{b'}(z)}{N_{b'}/N_b}\right\rfloor.
\]
In particular, a finer key determines every coarser key.  This nesting
is the reason for retaining dyadic grids even when $q$ is not dyadic.

\subsection*{Assigning the windows}

Take a complete ordered $M$-ary tree of height $d$, with leaves of height
$0$.  Write $P_1,\ldots,P_M$ for the children of an internal vertex $P$,
and $(P,i)$ for the edge from $P$ to $P_i$.  Let $K$ be the total number
of edges.  Order the edges in preorder: for each $i=1,\ldots,M$, list
$(P,i)$ and then all edges in the subtree rooted at $P_i$, in the same
recursive order.  An edge and its child subtree therefore form a
consecutive block in this list.

Assign an integer interval $W_e=[a_e,b_e]\cap\N$ to each edge, in
preorder, leaving exactly $g$ unused indices between consecutive
windows.  All edges leaving a vertex of height $h$ receive a window of
length $r_h$.  For an edge with a nonleaf child, we make its window at
least as long as the gap and the entire child subtree that follow it.
To specify these lengths from the leaves upward, write $\sigma_h$ for
the span, including gaps, of all windows in a subtree of height $h$.
At height $1$ there are $M$ windows and $M-1$ gaps, so we set
\[
 r_1=r_0,\qquad \sigma_1=Mr_1+(M-1)g,
\]
and, for $2\le h\le d$, choose
\[
 r_h=\max\{r_0,g+\sigma_{h-1}\}.
\]
At height $h\ge2$ there are $M$ blocks, each comprising one outgoing
window, a gap, and a child subtree.  Another $M-1$ gaps separate these
blocks, giving
\[
 \sigma_h=M(r_h+g+\sigma_{h-1})+(M-1)g.
\]
Choose the first starting index $n_0\ge1$ so that
$2q^{n_0}\le1/4$, and place all subsequent windows with the prescribed
lengths and gaps.  Attach to each edge $e$ the grid indexed by its window
endpoint $b_e$, and set
\[
 \mathcal N=\bigcup_e W_e.
\]

For $e=(P,i)$, let $b_e^*$ be the largest endpoint among $W_e$ and all
windows in the subtree rooted at $P_i$.  The recursive choice has the
following useful consequence.

\begin{lemma}[Span of an edge and its child subtree]\label{lem:window-span}
If $e=(P,i)$ leaves a vertex of height $h$, then
\begin{equation}\label{eq:window-span}
 b_e^*-a_e+1\le2r_h.
\end{equation}
\end{lemma}

\begin{proof}
If $h=1$, the child is a leaf, so $b_e^*=b_e$ and the span is $r_1$.
If $h\ge2$, preorder places the child's subtree immediately after
$W_e$, apart from one gap of length $g$.  The span is therefore
$r_h+g+\sigma_{h-1}$, which is at most $2r_h$ by the definition of
$r_h$.
\end{proof}

Figure~\ref{fig:window-block} shows the block appearing in the proof.
Thus the finest grid used anywhere in this block has its index at most
$2r_h-1$ beyond $a_e$.  This will control the number of grid boundaries
encountered as $t$ varies.  The gaps between windows serve a different
purpose: they control the bound on the density of centers at which a
translation from a later window can change an earlier key.

\begin{figure}[tb]
\centering
\begin{tikzpicture}[x=0.97cm,y=0.92cm,font=\small]
 \filldraw[fill=blue!12,draw=blue!65!black] (0,0) rectangle (4.8,1.05);
 \node at (2.4,.525) {$W_e$: $r_h$ indices};
 \draw[gray] (4.8,0) rectangle (5.4,1.05);
 \node at (5.1,.525) {$g$};
 \filldraw[fill=teal!12,draw=teal!60!black] (5.4,0) rectangle (9.3,1.05);
 \node[align=center] at (7.35,.525) {child-subtree windows\\and gaps: $\sigma_{h-1}$};
 \draw[gray] (9.3,0) rectangle (9.9,1.05);
 \node at (9.6,.525) {$g$};
 \filldraw[fill=black!5,draw=gray] (9.9,0) rectangle (12.3,1.05);
 \node[align=center] at (11.1,.525) {next sibling\\block};
 \node[above] at (0,1.1) {$a_e$};
 \node[above] at (4.8,1.1) {$b_e$};
 \node[above] at (9.3,1.1) {$b_e^*$};
 \draw[decorate,decoration={brace,mirror,amplitude=5pt},thick]
   (0,-.16)--(9.3,-.16);
 \node[below,align=center] at (4.65,-.42)
   {$b_e^*-a_e+1=r_h+g+\sigma_{h-1}\le2r_h$};
 \draw[-{Stealth[length=2mm]},gray] (-.2,-1.25)--(12.5,-1.25);
 \node[below] at (6.15,-1.25) {increasing indices in preorder};
\end{tikzpicture}
\caption{The block attached to an edge at height $h\ge2$, shown
schematically. The window length is at least the length of the following
gap and child-subtree span together. Thus the entire block has
at most twice the window length, regardless of where it occurs in the
tree. The next sibling block is outside this bound. At height $1$,
the child is a leaf and the block consists only of $W_e$.}
\label{fig:window-block}
\end{figure}


\subsection*{Preserving earlier keys}

Call $x\in\R$ \emph{stable} if, for every window $W_e$ other than the
first, writing $b'$ for the endpoint of its immediate predecessor,
\[
 (x,x+2q^{a_e}]\cap\Gamma_{b'}=\varnothing.
\]
Let $\mathcal S$ be the set of stable centers.  The left endpoint is
excluded because a translation to the right need not leave a grid cell
when it starts on that cell's left boundary; the right endpoint is
included because reaching the next boundary does change the key.

\begin{lemma}[Stability and earlier keys]\label{lem:stability}
The set $\mathcal S$ is measurable and $1$-periodic, and
\[
 \rho(\R\setminus\mathcal S)\le4Kc\,q^{g+1}.
\]
If $x\in\mathcal S$, $n\in W_e$, and $t\in[1,2]$, then every edge $f$
preceding $e$ in preorder satisfies
\begin{equation}\label{eq:stable-keys}
 J_{b_f}(x+tq^n)=J_{b_f}(x).
\end{equation}
\end{lemma}

\begin{proof}
For a fixed noninitial window, failure of its stability condition occurs
on the union of intervals
\[
 [\gamma-2q^{a_e},\gamma),\qquad \gamma\in\Gamma_{b'}.
\]
This is a measurable periodic set of density at most
\[
 2N_{b'}q^{a_e}
 \le4c\,q^{a_e-b'}
 =4c\,q^{g+1},
\]
where we used \eqref{eq:grid-size} and $a_e=b'+g+1$.  Summing over the
at most $K-1$ noninitial windows proves the stated density bound.

For the first window there are no earlier edges to consider.  Otherwise,
let $b'$ again be the endpoint of its predecessor, and suppose $x$ is
stable.  For $n\in W_e$ and $t\in[1,2]$, we have
$0<tq^n\le2q^{a_e}$.  Hence the segment from $x$ to $x+tq^n$ crosses
no boundary of $\Gamma_{b'}$, with a possible boundary at $x$ itself
assigned to the cell on its right.  It follows that
$J_{b'}(x+tq^n)=J_{b'}(x)$.  Every preceding edge $f$ has $b_f\le b'$,
so nesting gives \eqref{eq:stable-keys}.
\end{proof}

\subsection*{Separating the test points}

Earlier keys are now fixed at stable centers.  At the resolution attached
to the current edge, the center and its translated points instead have
distinct keys.  This second property holds at every center.

\begin{lemma}[Separation within one window]\label{lem:separation}
Fix an edge $e$, a center $x\in\R$, and a scale $t\in[1,2]$.  The points
\[
 \{x\}\ \cup\ \{x+tq^n:n\in W_e\}
\]
have pairwise distinct keys under $J_{b_e}$.  They also have pairwise
distinct keys under every $J_b$ with $b\ge b_e$.
\end{lemma}

\begin{proof}
Write $b=b_e$.  All the displayed points lie in
$[x,x+2q^{n_0}]$, an interval of length at most $1/4$.  Their pairwise
distances on the circle $\R/\Z$ therefore equal their ordinary real
distances.  The distance from $x$ to any translated point is
$tq^n\ge q^b$.  For $n<n'$ in $W_e$, the distance between the two
translated points is
\[
 t(q^n-q^{n'})=tq^n(1-q^{n'-n})\ge(1-q)q^b.
\]
On the other hand, \eqref{eq:grid-size} gives
\[
 \frac1{N_b}\le\frac{q^b}{c}
 =\frac{(1-q)q^b}{4}<(1-q)q^b.
\]
Two points with the same periodic key have circular distance less than
the cell width $1/N_b$.  Thus all the keys are distinct, including when
one of the points is a boundary point.  Distinctness persists at finer
resolutions because a finer key determines its coarser key.
\end{proof}

The random construction can consequently reuse all earlier decisions at
a stable center while consulting distinct entries for the points in the
current window.  We next use this distinction to obtain independent
opportunities for a hit.
